\section{Resumen}

Dos productos parecidos pueden verse iguales en el anaquel y traer datos distintos. Uno declara azúcar y fibra. Otro no dice nada de alérgenos. Otro no tiene precio. Compararlos a ojo mezcla lo que el producto afirma con lo que el registro simplemente no trae.

La pregunta es cómo convertir esa información, heterogénea e incompleta, en una comparación útil y repetible, sin inventar lo que los datos no traen. La persona indica qué quiere privilegiar, si sigue una dieta vegana o vegetariana y qué alérgenos necesita evitar. El orden usa solo la información disponible y muestra con qué datos se armó. No afirma que un producto sea mejor para toda la población.

El catálogo sale de Open Food Facts, una base pública de fichas de alimentos, en su corte de México. Si existe un precio observado, se toma de Open Prices o de Quién es Quién en los Precios. El análisis de esos datos fijó el diseño: no inventar lo que falta, comparar la nutrición entre productos parecidos, usar solo cinco nutrientes con una dirección clara, afinar el grado de procesamiento con el número de aditivos y dejar el precio fuera del orden. El cálculo es determinista: las mismas entradas dan siempre el mismo resultado. NutriMatch no utiliza aprendizaje supervisado porque el catálogo no contiene un target observado que permita entrenar y evaluar de manera científicamente válida un modelo de recomendación. Dicho de otra forma, no hay compras registradas que sirvan de respuesta correcta para entrenar un modelo.

El archivo de este trabajo es del \corte: 13\,093 productos de México y 273 columnas, sin códigos repetidos. De ellos, 5\,864 (44{,}79\,\%) tienen información suficiente para comparar nutrición y procesamiento antes de conocer a la persona. El precio real existe en 259 productos. Los otros 12\,834 no tienen precio. En un perfil de ejemplo, el orden reúne 113 productos, y 114 si cambian las prioridades.

Para revisar la parte nutricional se usó Nutri-Score como contraste externo, no como verdad del sistema. En 6\,035 productos, las dos medidas van en sentidos opuestos, que es lo esperado: la correlación de Spearman es $-0{,}475683$. Doce de trece casos revisados a mano coinciden. El caso restante difiere en el nombre, no en el puntaje.

La limitación principal es doble. Muchos registros no traen la información necesaria, y nadie observó la compra final. El sistema puede ordenar lo que está documentado. No puede decir si la persona elegirá el producto que quedó primero.
